\documentclass[11pt]{article}
\usepackage[a4paper,margin=28mm]{geometry}
\usepackage{amsmath,amssymb,amsthm}
\usepackage[colorlinks=true,linkcolor=blue,urlcolor=blue,citecolor=blue]{hyperref}
\newtheorem*{mainthm}{Theorem}
\newtheorem{lemma}{Lemma}
\newtheorem{proposition}[lemma]{Proposition}
\newcommand{\Q}{\mathbb Q}
\newcommand{\Z}{\mathbb Z}
\newcommand{\F}{\mathbb F}
\DeclareMathOperator{\End}{End}
\DeclareMathOperator{\Gal}{Gal}
\DeclareMathOperator{\GL}{GL}
\DeclareMathOperator{\PGL}{PGL}
\DeclareMathOperator{\Aut}{Aut}
\DeclareMathOperator{\Hom}{Hom}
\DeclareMathOperator{\tr}{tr}
\DeclareMathOperator{\spe}{sp}
\title{An obstruction to rational \(30\)-torsion on geometrically split
abelian surfaces}
\author{}
\date{21 October 2026}

\begin{document}
\maketitle

\begin{mainthm}
Let \(A/\Q\) be an abelian surface that is simple over \(\Q\) and
geometrically split. Then \(A(\Q)\) has no point of order \(30\).
\end{mainthm}

Here ``geometrically split'' means that \(A_{\overline{\Q}}\) is
isogenous to a product of elliptic curves. We impose no restriction on
the degree of this isogeny, on the integral endomorphism ring, or on the
polarization; the theorem does not require a principal polarization.
Each of the groups \([30]\), \([60]\) and \([2,30]\) contains a point of
order \(30\). Hence none of them occurs as the rational torsion subgroup
of a geometrically split genus-\(2\) Jacobian that is simple over
\(\Q\).

The hypothesis of simplicity over \(\Q\) is essential. A product of
elliptic curves over \(\Q\) with rational points of orders \(5\) and
\(6\) has a rational point of order \(30\), and the census contains
Jacobians split over \(\Q\) with each of the three torsion groups.
Simplicity is used only in Section~\ref{sec:nonisogenous}.
Section~\ref{sec:hypotheses} records the stronger statements that the
other sections prove for geometric squares.

\section{Inputs}\label{sec:inputs}

For a number field \(K\subset\overline{\Q}\) we write
\(G_K=\Gal(\overline{\Q}/K)\), and
\(\chi_p:G_{\Q}\to\F_p^\times\) denotes the mod-\(p\) cyclotomic
character. We use the following results.
\begin{enumerate}
\item \textbf{Elliptic torsion.} Derickx--Najman
\cite[Theorem 1.1]{DN} list the torsion groups of elliptic curves over
quartic number fields. The exponent of each group in the list is at
most \(24\), and none is divisible by \(30\). Hence no elliptic curve
over a number field of degree \(1\), \(2\) or \(4\) has a point of
order \(30\); for degrees \(1\) and \(2\), extend the ground field to
a field of degree \(4\). For degree \(1\) this is also Mazur's
Theorem~(8) \cite{Mazur}. For degree \(2\) it is also a special case
of Kenku--Momose \cite[Theorem (0.1)]{KM}, applied to
\((m,N)=(1,30)\): the curve \(X_1(30)\) has genus \(9\), and
\((1,30)\) is not one of the exceptions listed there.

\item \textbf{Good reduction.} Let \(K\) be a finite extension of
\(\Q_2\), let \(A/K\) be an abelian variety, and let \(\ell\) be an
odd prime. We use the following results of Serre--Tate \cite{ST}.
\begin{enumerate}
\item (Theorem 1 and its Corollary 2, p.~493.) \(A\) has good
reduction if and only if the inertia group of \(K\) acts trivially on
\(T_\ell(A)\). Good reduction is invariant under \(K\)-isogeny.
\item (Lemma 2, p.~495.) Suppose that \(A\) has good reduction, and let
\(B\) be the special fibre of its N\'eron model, an abelian variety
over the residue field \(k\). For every odd \(m\), reduction gives an
isomorphism \(A(\overline K)[m]\to B(\overline k)[m]\), compatible with
the action of \(\Gal(\overline K/K)\), which acts on the target
through \(\Gal(\overline k/k)\). In particular, \(A(K)[m]\) injects
into \(B(k)\), and \(T_\ell(A)\simeq T_\ell(B)\) equivariantly.
\item (Theorem 6(a), p.~503.) An abelian variety over \(K\) with
complex multiplication defined over \(K\), in the sense of
Serre--Tate, \S4, has potential good reduction. This applies to an
elliptic curve \(E/K\) such that \(\Q\otimes\End_K(E)\) is an
imaginary quadratic field.
\item (\S2, ``The case of a finite residue field'', p.~498.) If \(A\)
has potential good reduction, there is a finite totally ramified
extension \(K'/K\) such that \(A_{K'}\) has good reduction. This
consequence is stated explicitly in Clark--Xarles \cite[\S2.4]{CX},
and in Clark \cite[\S2.1]{Clark}.
\end{enumerate}

\item \textbf{Endomorphisms over finite fields.} Let \(B\) be an
abelian variety over a finite field \(k\), with Frobenius endomorphism
\(\pi\), and let \(f\) be the characteristic polynomial of \(\pi\). By
Tate \cite[Theorem 2(c)]{Tate}, the algebra \(\Q\otimes\End_k(B)\)
equals \(\Q[\pi]\) if and only if \(f\) has no multiple root.

\item \textbf{Weil's theorems.} Let \(B\) be an abelian surface over
\(\F_2\), with Frobenius endomorphism \(\pi\), and let \(\ell\) be an
odd prime. The characteristic polynomial of \(\pi\) on \(V_\ell(B)\)
has the form
\[
 P(T)=T^4+aT^3+bT^2+2aT+4,\qquad a,b\in\Z,
\]
its complex roots have absolute value \(\sqrt2\), and
\(\#B(\F_2)=P(1)\). The arithmetic Frobenius \(x\mapsto x^2\) of
\(\Gal(\overline{\F}_2/\F_2)\) acts on \(B(\overline{\F}_2)\), and
hence on \(T_\ell(B)\), as \(\pi\). An endomorphism of
\(B_{\overline{\F}_2}\) is defined over \(\F_2\) if and only if it
commutes with \(\pi\).

\item \textbf{Lattices.} Let \(G\) be a group and \(V\) a
finite-dimensional \(\Q_p\)-representation of \(G\). If
\(T_1,T_2\subset V\) are \(G\)-stable \(\Z_p\)-lattices, then
\((T_1/pT_1)^{\mathrm{ss}}\simeq(T_2/pT_2)^{\mathrm{ss}}\). Indeed,
the characteristic polynomial of each \(g\in G\) on \(T_i/pT_i\) is
the reduction of its characteristic polynomial on \(V\), and the
Brauer--Nesbitt theorem shows that a semisimple representation over a
field is determined by these polynomials. We apply this to isomorphic
representations by transporting a lattice through the isomorphism.
Composition factors are always counted with multiplicity. Rational
torsion is never treated as an isogeny invariant: only composition
factors are compared.
\end{enumerate}

Clark's Theorem~5 \cite{Clark} lists the possible orders of the
rational torsion subgroup of an abelian surface over \(\Q\) with
everywhere potentially good reduction; Theorem~1.4 of Clark--Xarles
\cite{CX} gives the same list for the larger class of surfaces with
everywhere anisotropic reduction. The list contains \(30\) and \(60\).
Thus, unlike the case of rational
\href{https://genus-2-torsion.github.io/proofs/qsimple-35-obstruction.html}{\(35\)-torsion}, it does not settle the CM
square case here. Section~\ref{sec:cm} replaces it with a local
argument at \(2\).

\section{Combining coprime torsion in an elliptic isogeny class}
\label{sec:combining}

\begin{lemma}\label{lem:combine}
Let \(m\) and \(n\) be coprime positive integers, and let \(K\) be a
number field. Suppose \(E_m\) and \(E_n\) are \(K\)-isogenous
elliptic curves over \(K\), with points of orders \(m\) and \(n\) in
\(E_m(K)\) and \(E_n(K)\), respectively. Then their \(K\)-isogeny
class contains an elliptic curve with a \(K\)-rational point of order
\(mn\).
\end{lemma}

\begin{proof}
Choose a \(K\)-isogeny \(f:E_m\to E_n\). Let \(C\subseteq\ker f\) be
the subgroup of points whose order divides a power of \(n\). It is
stable under \(\Gal(\overline K/K)\), because \(\ker f\) is. Hence
\(E'=E_m/C\) is an elliptic curve over \(K\), and \(f\) factors as
\[
 E_m\xrightarrow{\;a\;}E'=E_m/C\xrightarrow{\;b\;}E_n,
\]
with \(a\) and \(b\) defined over \(K\). The order of \(\ker a=C\) is
prime to \(m\), so \(a\) maps the given point \(P\) of order \(m\) to a
point of order \(m\). The order of \(\ker b\simeq\ker f/C\) is prime to
\(n\). Let \(Q\in E_n(K)\) be the given point of order \(n\), and let
\(\hat b\) be the dual isogeny, so that \(b\circ\hat b=[\deg b]\). If
\(\hat b(kQ)=0\), then \((\deg b)kQ=0\), hence \(kQ=0\). Thus
\(\hat b(Q)\in E'(K)\) has order \(n\), and
\(a(P)+\hat b(Q)\in E'(K)\) has order \(mn\). Since \(E'\) is
\(K\)-isogenous to \(E_m\), this proves the lemma.
\end{proof}

Applying the lemma first with \((m,n)=(2,3)\) and then with
\((m,n)=(6,5)\), we see: if a \(K\)-isogeny class of elliptic curves
contains curves with \(K\)-rational points of each of the orders
\(2\), \(3\) and \(5\), then it contains a curve with a
\(K\)-rational point of order \(30\).

For the rest of the proof, suppose that \(A(\Q)\) contains a point of
order \(30\), and choose nonzero points \(P_p\in A(\Q)[p]\) for
\(p\in\{2,3,5\}\). Exactly one of the following holds:
\begin{itemize}
\item \(A_{\overline{\Q}}\) is isogenous to \(E_1\times E_2\), where
\(E_1\) and \(E_2\) are not geometrically isogenous
(Section~\ref{sec:nonisogenous});
\item \(A_{\overline{\Q}}\) is isogenous to \(E^2\), where \(E\) has
complex multiplication (Section~\ref{sec:cm});
\item \(A_{\overline{\Q}}\) is isogenous to \(E^2\), where \(E\) has
no complex multiplication
(Sections~\ref{sec:endfield}--\ref{sec:completion}).
\end{itemize}

\section{Nonisogenous geometric factors}\label{sec:nonisogenous}

Suppose that \(A_{\overline{\Q}}\) is isogenous to \(E_1\times E_2\),
where \(E_1\) and \(E_2\) are not geometrically isogenous. Fix an
isogeny \(\psi:A_{\overline{\Q}}\to E_1\times E_2\). If
\(C\subset A_{\overline{\Q}}\) is an elliptic curve, then
\(\psi(C)\) is an elliptic curve isogenous to \(C\). Its two
projections cannot both be nonzero, since \(\Hom(E_1,E_2)=0\). Hence
\(\psi(C)\) equals \(E_1\times0\) or \(0\times E_2\), and \(C\) is the
identity component of \(\psi^{-1}(E_1\times0)\) or of
\(\psi^{-1}(0\times E_2)\). Thus \(A_{\overline{\Q}}\) has exactly two
elliptic subvarieties \(C_1\neq C_2\).

The group \(G_{\Q}\) permutes \(\{C_1,C_2\}\). If it fixed both, each
would be defined over \(\Q\), and \(A\) would contain an elliptic curve
over \(\Q\), contradicting simplicity over \(\Q\). This is the only
place where simplicity over \(\Q\) is used. Hence the stabilizer of
\(C_1\) is \(G_K\) for a quadratic field \(K\), and the nontrivial
automorphism \(\sigma\) of \(K\) interchanges the two subvarieties.
Write \(E=C_1\) and \(E^\sigma=C_2\); both are elliptic curves over
\(K\).

Put \(F=A_K/E\), with quotient map \(q:A_K\to F\). The restriction of
\(q\) to \(E^\sigma\) has finite kernel \(E\cap E^\sigma\) and nonzero
image, so it is a \(K\)-isogeny \(E^\sigma\to F\). Conjugating by
\(\sigma\), we obtain a \(K\)-isogeny \(E\to F^\sigma\).

Let \(p\in\{2,3,5\}\). If \(q(P_p)=0\), then \(P_p\in E(K)\) has
order \(p\). Otherwise \(q(P_p)\in F(K)\) has order \(p\), and its
conjugate \(\sigma(q(P_p))\in F^\sigma(K)\) also has order \(p\). In
both cases, a curve in the \(K\)-isogeny class of \(E\) has a
\(K\)-rational point of order \(p\). By Lemma~\ref{lem:combine}, this
class contains an elliptic curve over the quadratic field \(K\) with a
point of order \(30\). This contradicts input~1.

\section{The CM square case: a local obstruction at \(2\)}\label{sec:cm}

We prove a stronger local statement. It does not use simplicity over
\(\Q\) or the point of order \(2\).

\begin{proposition}\label{prop:cm}
Let \(A\) be an abelian surface over \(\Q_2\) such that
\(A_{\overline{\Q}_2}\) is isogenous to \(E^2\) for an elliptic curve
\(E\) with complex multiplication. Then \(A(\Q_2)\) has no point of
order \(15\).
\end{proposition}

If \(A/\Q\) is geometrically isogenous to the square of a CM elliptic
curve, then so is \(A_{\Q_2}\), and \(A(\Q)\subseteq A(\Q_2)\). Hence
the proposition excludes the CM square case of the theorem.

\begin{proof}
Let \(M\) be the imaginary quadratic field of complex multiplication
of \(E\). Then \(\End^0(A_{\overline{\Q}_2})\simeq M_2(M)\), with
center \(M\). Suppose that \(A(\Q_2)\) contains a point of order
\(15\).

\emph{Step 1: good reduction over a totally ramified extension.}
Choose a finite extension \(K_0/\Q_2\) such that \(E\) and all its
endomorphisms are defined over \(K_0\), and \(A_{K_0}\) is
\(K_0\)-isogenous to \(E^2\). By Serre--Tate, Theorem 6(a), \(E\) has
good reduction over a finite extension \(K_1/K_0\). Then \(E^2\) has
good reduction over \(K_1\), and so does the isogenous surface
\(A_{K_1}\), by Serre--Tate, Corollary 2 of Theorem 1. Thus \(A\) has
potential good reduction. By Serre--Tate, p.~498, there is a finite
totally ramified extension \(K'/\Q_2\) over which \(A\) has good
reduction. Since \(K'/\Q_2\) is totally ramified, the residue field of
\(K'\) is \(\F_2\). Let \(B/\F_2\) be the special fibre of the N\'eron
model of \(A_{K'}\).

\emph{Step 2: the number of points.} By input~2(b), \(A(K')[15]\)
injects into \(B(\F_2)\). Since \(A(\Q_2)\subseteq A(K')\), the group
\(B(\F_2)\) contains a point of order \(15\), so
\(15\mid\#B(\F_2)\). Let \(\pi\) be the Frobenius endomorphism of
\(B\), and \(P\) its characteristic polynomial, as in input~4. If
\(\alpha_1,\ldots,\alpha_4\) are the roots of \(P\), then
\[
 \#B(\F_2)=P(1)=\prod_{i=1}^4|1-\alpha_i|\le(1+\sqrt2)^4
 =17+12\sqrt2<34,
\]
so \(\#B(\F_2)\in\{15,30\}\).

\emph{Step 3: the only possible Weil polynomial.} Write
\(P(T)=T^4+aT^3+bT^2+2aT+4\) as in input~4. Then
\(P(T)=T^2\,Q(T+2/T)\) with \(Q(x)=x^2+ax+b-4\). If
\(\beta_1,\beta_2\) are the roots of \(Q\), then
\[
 P(T)=(T^2-\beta_1T+2)(T^2-\beta_2T+2),\qquad
 a=-(\beta_1+\beta_2),\qquad b=4+\beta_1\beta_2 .
\]
Each \(\beta_i\) equals \(\alpha+2/\alpha\) for a root \(\alpha\) of
\(P\). Since \(|\alpha|=\sqrt2\), we have \(2/\alpha=\overline\alpha\),
so \(\beta_i=2\operatorname{Re}\alpha\) is real and
\(|\beta_i|\le2\sqrt2\). Consequently \(|a|\le4\sqrt2\), so
\(|a|\le5\). Next, \(\beta_1\beta_2\le(\beta_1+\beta_2)^2/4\). Also
\((2\sqrt2-\beta_1)(2\sqrt2-\beta_2)\ge0\) and
\((2\sqrt2+\beta_1)(2\sqrt2+\beta_2)\ge0\), which give
\(\beta_1\beta_2\ge2\sqrt2|\beta_1+\beta_2|-8\). Hence
\begin{equation}\label{eq:bounds}
 2\sqrt2\,|a|-4\;\le\;b\;\le\;\frac{a^2}{4}+4.
\end{equation}
Moreover \(P(1)=5+3a+b\).

If \(P(1)=15\), then \(b=10-3a\). The upper bound in
\eqref{eq:bounds} gives \(a^2+12a-24\ge0\). Since \(a\) is an integer
with \(|a|\le5\), this forces \(a\ge2\). The lower bound then gives
\((3+2\sqrt2)a\le14\), so \(a\le14/(3+2\sqrt2)<2.41\). Hence \(a=2\)
and \(b=4\).

If \(P(1)=30\), then \(b=25-3a\). The upper bound gives
\(a^2+12a-84\ge0\), which forces \(a\ge5\). The lower bound gives
\((3+2\sqrt2)a\le29\), so \(a\le29/(3+2\sqrt2)<4.98\). This is a
contradiction. Therefore
\begin{equation}\label{eq:weil}
 P(T)=T^4+2T^3+4T^2+4T+4=(T^2+2)(T^2+2T+2),\qquad
 \tr(\pi\mid V_\ell(B))=-2.
\end{equation}

\emph{Step 4: the endomorphism algebra of \(B\).} The two factors in
\eqref{eq:weil} are distinct irreducible quadratic polynomials, so
\(P\) has no multiple root. By input~3,
\(\End^0_{\F_2}(B)=\Q[\pi]\). The minimal polynomial of \(\pi\)
divides \(P\) and vanishes at every root of \(P\), so it equals \(P\).
Therefore
\begin{equation}\label{eq:endB}
 \End^0_{\F_2}(B)=\Q[\pi]\simeq\Q[T]/(T^2+2)\times\Q[T]/(T^2+2T+2)
 \simeq\Q(\sqrt{-2})\times\Q(i).
\end{equation}

\emph{Step 5: specialization of the CM centre.} Fix an odd prime
\(\ell\). Since \(A_{K'}\) has good reduction, the inertia group of
\(K'\) acts trivially on \(T_\ell(A)\). An element \(\tau\) of
\(\Gal(\overline{\Q}_2/K')\) sends an endomorphism \(\varphi\) of
\(A_{\overline{\Q}_2}\) to \(\tau\varphi\tau^{-1}\), and on
\(T_\ell(A)\) this is conjugation by the action of \(\tau\).
Endomorphisms act faithfully on \(T_\ell(A)\). Hence inertia fixes
every geometric endomorphism, and all geometric endomorphisms of \(A\)
are defined over a finite unramified extension \(K''/K'\). The N\'eron
model of \(A_{K''}\) is the base change of that of \(A_{K'}\), so its
special fibre is the base change of \(B\) to the residue field of
\(K''\). By the N\'eron mapping property, each endomorphism of
\(A_{K''}\) extends to the N\'eron model and restricts to an
endomorphism of the special fibre. This gives a unital ring
homomorphism
\[
 \spe:\End^0(A_{\overline{\Q}_2})\longrightarrow
 \End^0(B_{\overline{\F}_2}).
\]
Under the isomorphism \(T_\ell(A)\simeq T_\ell(B)\) of input~2(b),
\(\spe(\varphi)\) acts on \(T_\ell(B)\) as \(\varphi\) acts on
\(T_\ell(A)\). Hence \(\spe\) is injective. Since that isomorphism is
Galois-equivariant, \(\spe\) is equivariant for the action of
\(\Gal(\overline{\Q}_2/K')\), which acts on
\(\End^0(B_{\overline{\F}_2})\) through \(\Gal(\overline{\F}_2/\F_2)\).
By input~4, the arithmetic Frobenius acts there as
\(\psi\mapsto\pi\psi\pi^{-1}\).

The centre \(M\) of \(\End^0(A_{\overline{\Q}_2})\) is stable under
\(\Gal(\overline{\Q}_2/K')\). This group acts on \(M\) through
\(\Gal(\overline{\F}_2/\F_2)\), since inertia acts trivially, and by
field automorphisms. Hence a Frobenius element acts on \(M\) either
trivially or by complex conjugation \(m\mapsto\overline m\).
Accordingly, either \(\pi\spe(m)\pi^{-1}=\spe(m)\) for all \(m\in M\),
or \(\pi\spe(m)\pi^{-1}=\spe(\overline m)\) for all \(m\in M\). We need
not know which case occurs: the CM field need not be contained in
\(K'\).

\emph{Step 6: trivial action.} Suppose \(\pi\) commutes with
\(\spe(M)\). By input~4, \(\spe(M)\) is contained in
\(\End^0_{\F_2}(B)\), which is the algebra \eqref{eq:endB}. Composing
the unital embedding \(M\to\Q(\sqrt{-2})\times\Q(i)\) with the two
projections gives unital ring homomorphisms from the field \(M\) to
\(\Q(\sqrt{-2})\) and to \(\Q(i)\). They are injective, and all three
fields have degree \(2\), so \(M\simeq\Q(\sqrt{-2})\) and
\(M\simeq\Q(i)\). This is impossible.

\emph{Step 7: nontrivial action.} Suppose
\(\pi\spe(m)\pi^{-1}=\spe(\overline m)\) for all \(m\in M\). Put
\(V=V_\ell(B)\otimes_{\Q_\ell}\overline{\Q}_\ell\), and let
\(\iota_1,\iota_2\) be the two embeddings of \(M\) into
\(\overline{\Q}_\ell\), with \(\iota_2(m)=\iota_1(\overline m)\).
Since \(M\otimes\overline{\Q}_\ell\simeq\overline{\Q}_\ell\times
\overline{\Q}_\ell\), we have \(V=V_1\oplus V_2\), where \(\spe(m)\)
acts on \(V_j\) as \(\iota_j(m)\). If \(v\in V_1\), then
\[
 \spe(m)\,\pi v=\pi\,\spe(\overline m)\,v=\iota_1(\overline m)\,\pi v
 =\iota_2(m)\,\pi v,
\]
so \(\pi V_1\subseteq V_2\), and similarly \(\pi V_2\subseteq V_1\).
With respect to the decomposition \(V=V_1\oplus V_2\), the matrix of
\(\pi\) has zero diagonal blocks, so \(\tr(\pi\mid V)=0\). This
contradicts \eqref{eq:weil}.

Both cases are impossible, which proves the proposition.
\end{proof}

The reduction \(B\) may have more endomorphisms than \(A\); only the
injective specialization of the centre \(M\) is used. We do not enlarge
\(K'\) to make \(M\) defined over it, because such an enlargement need
not remain totally ramified.

\section{The non-CM square case: the endomorphism field}
\label{sec:endfield}

In this section and the next four, \(A_{\overline{\Q}}\) is isogenous
to \(E_0^2\) for an elliptic curve \(E_0\) without complex
multiplication. Then
\[
 D=\End^0(A_{\overline{\Q}})\simeq M_2(\End^0(E_0))=M_2(\Q).
\]
The group \(G_{\Q}\) acts on the finitely generated \(\Z\)-module
\(\End(A_{\overline{\Q}})\) by \(g\cdot\varphi=g\varphi g^{-1}\), and
each endomorphism is defined over a number field. Let \(L\) be the
fixed field of the kernel of this action, that is, the smallest field
over which all geometric endomorphisms of \(A\) are defined. It is a
finite Galois extension of \(\Q\). Put \(\Gamma=\Gal(L/\Q)\). It acts
faithfully on \(D\) by \(\Q\)-algebra automorphisms. By
Skolem--Noether, \(\Aut_{\Q}(M_2(\Q))=\PGL_2(\Q)\), so \(\Gamma\) is a
finite subgroup of \(\PGL_2(\Q)\). For \(\gamma\in\Gamma\) we choose a
representative \(S_\gamma\in\GL_2(\Q)\), so that \(\gamma\) acts on
\(D\) as \(x\mapsto S_\gamma xS_\gamma^{-1}\).

A finite subgroup of \(\PGL_2(\mathbb R)\) is cyclic or dihedral.
Indeed, it acts by isometries on the hyperbolic plane, it fixes a point
there, and its action on the tangent plane at that point embeds it in
\(O(2)\). Let an element of \(\PGL_2(\Q)\) have order \(n>2\), and let
\(S\) be a representative. Then the ratio of the two eigenvalues of
\(S\) is a primitive \(n\)-th root of unity \(\zeta\), and
\[
 \frac{(\tr S)^2}{\det S}=2+\zeta+\zeta^{-1}\in\Q .
\]
Hence \(n\in\{3,4,6\}\). It follows that
\[
 \Gamma\in\{C_1,C_2,C_3,C_4,C_6,D_4,D_6,D_8,D_{12}\},
\]
where \(C_n\) is cyclic of order \(n\) and the subscript on a dihedral
group denotes its order; thus \(D_4=C_2\times C_2\). When \(\Gamma\)
is \(C_n\) or \(D_{2n}\) with \(n\ge3\), its rotation subgroup is the
unique cyclic subgroup \(\langle r\rangle\) of order \(n\). It is
normal, and \(\gamma r\gamma^{-1}\in\{r,r^{-1}\}\) for every
\(\gamma\in\Gamma\).

\paragraph{Trace-one representatives.} Let \(r\in\Gamma\) have order
\(n\ge3\), and let \(S\) be a representative. Then \(\tr S\neq0\):
otherwise \(S^2=-\det S\) is scalar by Cayley--Hamilton, and \(r\) would
have order at most \(2\). Put \(R=S/\tr S\). Its characteristic
polynomial is \(X^2-X+d\), where
\(1/d=(\tr R)^2/\det R=2+\zeta+\zeta^{-1}\). Thus
\[
 X^2-X+1\ (n=3),\qquad X^2-X+\tfrac12\ (n=4),\qquad
 X^2-X+\tfrac13\ (n=6),
\]
with discriminants \(-3\), \(-1\) and \(-1/3\), respectively. Since
\(R(1-R)=d\), the trace-one representative of \(r^{-1}\) is \(1-R\).

\paragraph{Elliptic subvarieties.} Let \(C\subset A_{\overline{\Q}}\)
be an elliptic curve. By Poincar\'e's reducibility theorem there is an
elliptic curve \(C'\subset A_{\overline{\Q}}\), a complement of \(C\),
such that the addition map \(\mu:C\times C'\to A_{\overline{\Q}}\) is
an isogeny. Choose an isogeny \(\psi:A_{\overline{\Q}}\to C\times C'\)
with \(\psi\circ\mu=[N]\) for some integer \(N\ge1\). Composing
\(\psi\) with the first projection and the inclusion of \(C\) gives an
endomorphism \(\varepsilon_C\) of \(A_{\overline{\Q}}\) with image
\(C\), whose restriction to \(C\) is multiplication by \(N\). Since
\(\varepsilon_C\) is defined over \(L\), so is its image \(C\). Thus
every elliptic curve in \(A_{\overline{\Q}}\) is defined over \(L\).
Each of them is isogenous to \(E_0\): composing its inclusion with an
isogeny \(A_{\overline{\Q}}\to E_0^2\) and a suitable projection gives
a nonzero homomorphism to \(E_0\).

\section{Polarization duality and an irreducible projective factor}
\label{sec:cyclotomic}

Let \(p\) be a prime, \(T=T_p(A)\) and \(V=V_p(A)\). The algebra
\(D\otimes\Q_p\simeq M_2(\Q_p)\) acts faithfully on \(V\), and
\(G_{\Q}\) acts compatibly: \(g(xv)=(g\cdot x)(gv)\) for
\(x\in D\otimes\Q_p\) and \(v\in V\), where \(G_{\Q}\) acts on \(D\)
through \(\Gamma\).

\begin{lemma}[Polarization duality]\label{lem:duality}
Let \(T''\subset V\) be a \(G_{\Q}\)-stable lattice, let
\(H\subseteq G_{\Q}\) be a subgroup, and let \(Y\) be the
semisimplification of \(T''/pT''\) as an \(H\)-module. Then
\(Y\simeq Y^*\otimes\chi_p|_H\).
\end{lemma}

\begin{proof}
Every abelian variety over \(\Q\) has a polarization over \(\Q\).
Combined with the Weil pairing, it gives a nondegenerate alternating
\(G_{\Q}\)-equivariant pairing \(V\times V\to\Q_p(1)\), hence an
isomorphism \(V\simeq V^*\otimes\Q_p(1)\). The lattice
\(\Hom_{\Z_p}(T'',\Z_p)\otimes\Z_p(1)\) is \(G_{\Q}\)-stable, with
reduction \((T''/pT'')^*\otimes\chi_p\). Transport it into \(V\)
through this isomorphism. By input~5 over \(H\), its residual
semisimplification is \(Y\). On the other hand, dualizing and twisting
are exact, so it is also \(Y^*\otimes\chi_p|_H\).
\end{proof}

The polarization is used only on the rational Tate module, so its
degree may be divisible by \(p\).

\begin{lemma}\label{lem:cyclotomic}
Let \(p\) be an odd prime. Suppose \(\Gamma\) preserves a maximal
\(\Z_p\)-order \(\mathcal R\) of \(D\otimes\Q_p\), and the induced
projective action of \(\Gamma\) on the natural two-dimensional residual
module of \(\mathcal R\) is irreducible over \(\F_p\). If
\(A(\Q)[p]\neq0\), then \(\Q(\zeta_p)\subseteq L\).
\end{lemma}

To make the hypothesis precise, write
\(\mathcal R=\End_{\Z_p}(\Lambda)\) for a lattice
\(\Lambda\subset\Q_p^2\). An element \(\gamma\in\Gamma\) preserves
\(\mathcal R\) exactly when \(S_\gamma\Lambda\) is homothetic to
\(\Lambda\). We may then rescale \(S_\gamma\) so that
\(S_\gamma\in\GL(\Lambda)\). It is unique up to \(\Z_p^\times\), and
its reduction gives a homomorphism \(\Gamma\to\PGL(\Lambda/p\Lambda)\).
The hypothesis is that no line in \(U=\Lambda/p\Lambda\) is stable
under this projective action.

\begin{proof}
Put \(T'=\mathcal R T\). It is a full lattice in \(V\), because
\(\mathcal R\) is finite over \(\Z_p\). It is \(G_{\Q}\)-stable,
because \(T\) and \(\mathcal R\) are. Choosing a basis of \(\Lambda\)
identifies \(\mathcal R\) with \(M_2(\Z_p)\). The matrix unit
\(e_{21}\) maps \(e_{11}T'\) isomorphically onto \(e_{22}T'\), with
inverse \(e_{12}\). This gives
\[
 T'\simeq\Lambda\otimes_{\Z_p}W_0,\qquad W_0=e_{11}T',
\]
where \(W_0\) is free of rank \(2\) and \(\mathcal R\) acts on the
first factor.

Put \(\overline V=T'/pT'=U\otimes_{\F_p}\overline W_0\). By input~5,
\(\overline V\) and \(T/pT=A[p]\) have the same composition factors.
Since \(A(\Q)[p]\neq0\), the module \(\overline V\) has a trivial
composition factor.

For \(g\in G_{\Q}\) with image \(\gamma\in\Gamma\), and
\(x\in\mathcal R\), we have
\(g\circ\overline x\circ g^{-1}=\overline{S_\gamma xS_\gamma^{-1}}\) on
\(\overline V\). Pull back the central extension
\[
 1\longrightarrow\F_p^\times\longrightarrow\GL_2(\F_p)\longrightarrow
 \PGL_2(\F_p)\longrightarrow1
\]
along the projective action of \(G_{\Q}\) on \(U\). The result is the
group \(\widetilde G\) of pairs \((g,\theta)\), with \(g\in G_{\Q}\)
and \(\theta\in\GL(U)\) a lift of the projective class of
\(\overline{S_\gamma}\). For such a pair,
\(g(\overline x\otimes1)g^{-1}=\theta\overline x\theta^{-1}\otimes1\)
for \(x\in\mathcal R\), because scalars cancel. Hence the operator
\((\theta^{-1}\otimes1)\circ g\) on \(\overline V\) commutes with
\(M_2(\F_p)\otimes1\). The commutant of \(M_2(\F_p)\otimes1\) in
\(\End(U\otimes\overline W_0)\) is \(1\otimes\End(\overline W_0)\).
Hence \((\theta^{-1}\otimes1)\circ g=1\otimes w(g,\theta)\) for a
unique \(w(g,\theta)\in\GL(\overline W_0)\). By uniqueness, \(w\) is a
homomorphism. Let \(W\) denote \(\overline W_0\) with this action of
\(\widetilde G\), and let \(\widetilde G\) act on \(U\) through
\(\theta\). Then
\[
 \overline V\simeq U\otimes_{\F_p}W
\]
as \(\widetilde G\)-modules, where \(\widetilde G\) acts on
\(\overline V\) through its projection to \(G_{\Q}\). Since this
projection is surjective, a composition series of \(\overline V\) as a
\(G_{\Q}\)-module is one as a \(\widetilde G\)-module, so
\(\overline V\) has a trivial \(\widetilde G\)-composition factor.

The module \(U\) is irreducible by hypothesis. Since \(G_L\) acts
trivially on \(D\), the pairs \((h,1)\) with \(h\in G_L\) lie in
\(\widetilde G\). They form a normal subgroup
\[
 \widetilde H=\{(h,1):h\in G_L\}\subseteq\widetilde G,
\]
which acts trivially on \(U\).

If \(W\) had a \(\widetilde G\)-stable line \(W_1\), then
\(\overline V\) would have the filtration
\(0\subset U\otimes W_1\subset U\otimes W\). Both graded pieces are
twists of \(U\) by characters, hence irreducible of dimension \(2\).
By the Jordan--H\"older theorem, \(\overline V\) would have no trivial
composition factor. Thus \(W\) is irreducible.

The restriction \(W|_{\widetilde H}\) is semisimple. Indeed, choose a
simple \(\widetilde H\)-submodule of \(W\). Its
\(\widetilde G\)-translates are simple \(\widetilde H\)-submodules,
because \(\widetilde H\) is normal. Their sum is a nonzero
\(\widetilde G\)-submodule of \(W\), hence all of \(W\).

Therefore \(\overline V|_{\widetilde H}\simeq W|_{\widetilde H}\oplus
W|_{\widetilde H}\) is semisimple. It has a trivial composition factor,
so \(\overline V^{\widetilde H}\neq0\). Since
\(\overline V^{\widetilde H}=U\otimes W^{\widetilde H}\), also
\(W^{\widetilde H}\neq0\). The subspace \(W^{\widetilde H}\) is
\(\widetilde G\)-stable, because \(\widetilde H\) is normal, and \(W\)
is irreducible, so \(W^{\widetilde H}=W\). Hence \(G_L\) acts trivially
on \(\overline V\).

Apply Lemma~\ref{lem:duality} to \(T''=T'\) and \(H=G_L\). Then \(Y\)
is four copies of the trivial character, while \(Y^*\otimes\chi_p\)
is four copies of \(\chi_p|_{G_L}\). Hence \(\chi_p|_{G_L}=1\), which
is equivalent to \(\Q(\zeta_p)\subseteq L\).
\end{proof}

\section{Rotations of order \(3\) or \(6\)}\label{sec:rotations}

Suppose that
\[
 \Gamma\in\{C_3,C_6,D_6,D_{12}\},
\]
and let \(r\) generate the rotation subgroup, of order
\(n\in\{3,6\}\). Let \(R\) be its trace-one representative. Its
characteristic polynomial is \(X^2-X+1\) or \(X^2-X+\frac13\), with
discriminant \(-3\) or \(-\frac13\), respectively. Modulo \(5\) these
discriminants are \(2\) and \(3\), and neither is a square. Hence the
reduction of the characteristic polynomial modulo \(5\) is irreducible
over \(\F_5\). Thus \(\Q_5[R]\) is the unramified quadratic extension
of \(\Q_5\), and \(\Z_5[R]\) is its ring of integers.

An \(R\)-stable \(\Z_5\)-lattice \(\Lambda\subset\Q_5^2\) is a finitely
generated torsion-free module over the discrete valuation ring
\(\Z_5[R]\), and \(\Q_5^2\) has dimension \(1\) over \(\Q_5[R]\).
Hence \(\Lambda\) is free of rank \(1\) over \(\Z_5[R]\). Since the
uniformizer of \(\Z_5[R]\) is \(5\), any two such lattices are
homothetic under \(\Q_5^\times\). Now suppose \(r\) fixes the
homothety class of a lattice \(\Lambda\), so \(R\Lambda=c\Lambda\) for
some \(c\in\Q_5^\times\). Comparing indices gives
\(v_5(\det R)=2v_5(c)\). Since \(\det R\in\{1,\frac13\}\) is a
\(5\)-adic unit, \(c\) is a unit and \(R\Lambda=\Lambda\). Hence \(r\)
fixes exactly one homothety class \([\Lambda_0]\).

Let \(\gamma\in\Gamma\). Since \(\gamma^{-1}r\gamma=r^{\pm1}\), we have
\(S_\gamma^{-1}RS_\gamma=c'R^{\pm1}\) for a scalar \(c'\). As
\(R^{\pm1}\Lambda_0=\Lambda_0\), we get
\(R(S_\gamma\Lambda_0)=c'S_\gamma\Lambda_0\). Thus the class of
\(S_\gamma\Lambda_0\) is fixed by \(r\), hence equals \([\Lambda_0]\).
It follows that
\[
 \mathcal R=\End_{\Z_5}(\Lambda_0)
\]
is a \(\Gamma\)-stable maximal order. The projective action of \(r\) on
\(\Lambda_0/5\Lambda_0\) is given by the reduction of
\(R\in\GL(\Lambda_0)\). Its characteristic polynomial is irreducible
over \(\F_5\), so it has no eigenvector. Hence no line in
\(\Lambda_0/5\Lambda_0\) is stable, and the residual projective action
is irreducible.

Since \(A(\Q)[5]\neq0\), Lemma~\ref{lem:cyclotomic} gives
\(\Q(\zeta_5)\subseteq L\). Then \(\Gamma\) has a quotient
\(\Gal(\Q(\zeta_5)/\Q)\simeq C_4\). But \(C_3\) and \(C_6\) have order
prime to \(4\), and the abelianizations of \(D_6\) and \(D_{12}\) are
\(C_2\) and \(C_2\times C_2\). None of these groups has a cyclic
quotient of order \(4\), a contradiction.

\section{The dihedral group of order \(8\)}\label{sec:dihedral}

Suppose \(\Gamma=D_8\), with rotation subgroup \(\langle r\rangle\) of
order \(4\). Let \(R\) be the trace-one representative of \(r\). Its
characteristic polynomial \(X^2-X+\frac12\) has discriminant \(-1\),
which is congruent to \(2^2\) modulo \(5\). By Hensel's lemma, its
roots \(\lambda_1,\lambda_2\) lie in \(\Z_5\). They are units, since
\(\lambda_1\lambda_2=\frac12\), and \(\lambda_1-\lambda_2=\pm\sqrt{-1}\)
is a unit. Their ratio \(\zeta=\lambda_1/\lambda_2\) satisfies
\(\zeta+\zeta^{-1}=1/d-2=0\), so \(\zeta^2=-1\), and \(\zeta\) is a
primitive fourth root of unity. The ring
\[
 \mathcal B=\Z_5[R]=\Z_5e_1\oplus\Z_5e_2,\qquad
 e_1=\frac{R-\lambda_2}{\lambda_1-\lambda_2},\quad
 e_2=\frac{R-\lambda_1}{\lambda_2-\lambda_1},
\]
is a split quadratic order with orthogonal idempotents \(e_1,e_2\), and
\(R=\lambda_1e_1+\lambda_2e_2\). Since \(R\) is not scalar, each
\(e_i\) is a rank-one idempotent of \(M_2(\Q_5)\).

The rotations act by conjugation by powers of \(R\), so they fix
\(R\), \(e_1\) and \(e_2\). A reflection \(s\) satisfies
\(srs^{-1}=r^{-1}\). Hence \(S_sRS_s^{-1}\) represents \(r^{-1}\) and
has trace \(1\), so it is the trace-one representative \(1-R\) of
\(r^{-1}\). Since \(\lambda_1+\lambda_2=1\), it sends \(e_1\) to
\((\lambda_1-R)/(\lambda_1-\lambda_2)=e_2\). Hence \(\mathcal B\) is
\(\Gamma\)-stable: rotations fix \(e_1\) and \(e_2\), and reflections
interchange them. This holds even though \(\Gamma\) need not preserve
any maximal order.

Let \(K=L^{\langle r\rangle}\), a quadratic field. Put
\(T'=\mathcal BT_5(A)\). As in Section~\ref{sec:cyclotomic}, it is a
\(G_{\Q}\)-stable full lattice in \(V=V_5(A)\). We have
\[
 T'=T_1\oplus T_2,\qquad T_i=e_iT'.
\]
Each \(T_i\) is a direct summand of \(T'\), hence free. The
four-dimensional \(M_2(\Q_5)\)-module \(V\) is the sum of two copies of
the standard module, so \(V_i=e_iV\) has dimension \(2\) and \(T_i\)
has rank \(2\). The group \(G_K\) preserves each \(T_i\). Every element
\(\tau\in G_{\Q}\setminus G_K\) satisfies \(\tau T_1=T_2\) and
\(\tau T_2=T_1\).

The space \(e_1(D\otimes\Q_5)e_2\) is one-dimensional. Let \(X\) span
it. The group \(G_K\) acts on \(D\otimes\Q_5\) through
\(\langle r\rangle\), fixing \(e_1\) and \(e_2\), so
\[
 g\cdot X=\eta(g)X\qquad(g\in G_K)
\]
for a character \(\eta:G_K\to\Z_5^\times\), trivial on \(G_L\). Since
\(RXR^{-1}=(\lambda_1/\lambda_2)X\), its value on a lift of \(r\) is
\(\zeta\). Thus \(\eta\) has exact order \(4\). Its reduction
\(\overline\eta\) also has exact order \(4\), because \(\zeta^2=-1\)
and \(-1\not\equiv1\) modulo \(5\).

As a nonzero element of \(e_1M_2(\Q_5)e_2\), the matrix \(X\) maps the
line \(e_2\Q_5^2\) isomorphically onto \(e_1\Q_5^2\). Since \(V\) is a
sum of two copies of \(\Q_5^2\), the map \(X:V_2\to V_1\) is an
isomorphism. For \(g\in G_K\) and \(v\in V_2\) we have
\(g(Xv)=\eta(g)X(gv)\). Hence \(X\) is a \(G_K\)-isomorphism
\(\eta\otimes V_2\to V_1\). Put \(\overline T_i=T_i/5T_i\). Input~5,
applied to the lattices \(T_1\) and \(X(T_2)\) of \(V_1\), gives
\begin{equation}\label{eq:twist}
 \overline T_1^{\,\mathrm{ss}}\simeq
 \overline\eta\otimes\overline T_2^{\,\mathrm{ss}}
\end{equation}
as \(G_K\)-modules. This step does not require \(X\) to carry \(T_2\)
onto \(T_1\).

By input~5, \(T'/5T'=\overline T_1\oplus\overline T_2\) has the same
composition factors as \(A[5]\), so it has a trivial
\(G_K\)-composition factor. Hence one of \(\overline T_1,\overline
T_2\) has one. An element \(\tau\in G_{\Q}\setminus G_K\) maps
\(\overline T_1\) onto \(\overline T_2\) and intertwines the action of
\(g\in G_K\) on \(\overline T_1\) with that of \(\tau g\tau^{-1}\) on
\(\overline T_2\). Conjugation preserves trivial composition factors,
so both \(\overline T_1\) and \(\overline T_2\) have one. By
\eqref{eq:twist}, \(\overline T_1^{\,\mathrm{ss}}\) contains both
\(\mathbf 1\) and \(\overline\eta\otimes\mathbf 1=\overline\eta\).
Since \(\overline\eta\neq\mathbf 1\) and \(\dim\overline T_1=2\),
\[
 \overline T_1^{\,\mathrm{ss}}\simeq\mathbf 1\oplus\overline\eta,
 \qquad
 \overline T_2^{\,\mathrm{ss}}\simeq\mathbf 1\oplus\overline\eta^{-1}.
\]
Since \(\overline\eta\) is trivial on \(G_L\), the restriction of
\(T'/5T'\) to \(G_L\) has four trivial composition factors.
Lemma~\ref{lem:duality}, with \(T''=T'\) and \(H=G_L\), gives
\(\chi_5|_{G_L}=1\), that is, \(\Q(\zeta_5)\subseteq L\). Then \(D_8\)
would have a cyclic quotient of order \(4\). This is impossible, since
its abelianization is \(C_2\times C_2\).

Alternatively, apply Lemma~\ref{lem:duality} with \(H=G_K\). The
\(G_K\)-composition factors of \(T'/5T'\) are
\(\mathbf 1,\mathbf 1,\overline\eta,\overline\eta^{-1}\), and the
trivial character is the only one of multiplicity at least \(2\). The
dual twist consists of
\(\chi_5,\chi_5,\chi_5\overline\eta^{-1},\chi_5\overline\eta\), in
which \(\chi_5|_{G_K}\) has multiplicity at least \(2\). Hence
\(\chi_5|_{G_K}=1\), which would put the quartic field \(\Q(\zeta_5)\)
inside the quadratic field \(K\).

\section{Completion of the proof}\label{sec:completion}

Sections~\ref{sec:rotations} and~\ref{sec:dihedral} leave
\[
 \Gamma\in\{C_1,C_2,C_4,D_4\},
\]
so \([L:\Q]\in\{1,2,4\}\). By Section~\ref{sec:endfield}, every
elliptic curve in \(A_{\overline{\Q}}\) is defined over \(L\). Choose
one, \(E\), and put \(F=A_L/E\), with quotient map \(q:A_L\to F\).

The elliptic curves \(E\) and \(F\) are \(L\)-isogenous. Indeed, let
\(E'\) be a complement of \(E\) as in Section~\ref{sec:endfield}, so
that \(E\cap E'\) is finite; it is defined over \(L\). Both \(E\) and
\(E'\) are isogenous to \(E_0\), so there is a nonzero homomorphism
\(h:E\to E'\) over \(\overline{\Q}\). Let \(\varepsilon_E\) be the
endomorphism of Section~\ref{sec:endfield}, with image \(E\) and
restriction \([N]\) to \(E\). Then \(\varphi=h\circ\varepsilon_E\),
viewed as an endomorphism of \(A_{\overline{\Q}}\), is defined over
\(L\), and \(\varphi(E)=E'\). The composite
\[
 E\xrightarrow{\;\varphi|_E\;}E'\xrightarrow{\;q|_{E'}\;}F
\]
is defined over \(L\). Both maps are isogenies: the first is \(Nh\),
and the second has finite kernel \(E\cap E'\).

For each \(p\in\{2,3,5\}\), the exact sequence
\[
 0\longrightarrow E\longrightarrow A_L\xrightarrow{\;q\;}F
 \longrightarrow0
\]
shows that \(P_p\) lies in \(E(L)\) or maps to a point of order \(p\)
in \(F(L)\). Thus the common \(L\)-isogeny class of \(E\) and \(F\)
contains curves with \(L\)-rational points of orders \(2\), \(3\) and
\(5\). By Lemma~\ref{lem:combine}, it contains an elliptic curve over
\(L\) with a point of order \(30\). Since \([L:\Q]\in\{1,2,4\}\), this
contradicts input~1.

All three geometric types of Section~\ref{sec:combining} lead to a
contradiction, which proves the theorem.
\hfill\(\square\)

\section{Where the hypotheses are used}\label{sec:hypotheses}

Simplicity over \(\Q\) is used only in
Section~\ref{sec:nonisogenous}, where it forces Galois to interchange
the two geometric elliptic subvarieties. The other sections prove two
stronger statements:
\begin{itemize}
\item If \(A/\Q_2\) is an abelian surface that is geometrically
isogenous to the square of a CM elliptic curve, then \(A(\Q_2)\) has no
point of order \(15\) (Proposition~\ref{prop:cm}). In particular, such
an \(A/\Q\) has no rational point of order \(15\).
\item If \(A/\Q\) is an abelian surface that is geometrically isogenous
to the square of an elliptic curve without complex multiplication, then
\(A(\Q)\) has no point of order \(30\)
(Sections~\ref{sec:endfield}--\ref{sec:completion}), whether or not
\(A\) is simple over \(\Q\).
\end{itemize}
Hence an abelian surface over \(\Q\) with a rational point of order
\(30\) that is geometrically split is geometrically isogenous to a
product of two elliptic curves that are not geometrically isogenous,
and it is not simple over \(\Q\). In Section~\ref{sec:nonisogenous} the
points of orders \(2\), \(3\) and \(5\) enter only through
Lemma~\ref{lem:combine}. In Section~\ref{sec:cm} the points of orders
\(3\) and \(5\) together give the contradiction, and the point of order
\(2\) is not used. In the non-CM square case, the point of order \(5\)
drives Sections~\ref{sec:rotations} and~\ref{sec:dihedral}, and the
points of orders \(2\) and \(3\) enter only through
Lemma~\ref{lem:combine} in Section~\ref{sec:completion}.

No step uses the degree of a geometric splitting isogeny, the
maximality of the integral endomorphism ring, or a principal
polarization. No numerical computation is required: the only finite
check, the determination of the Weil polynomial in
Section~\ref{sec:cm}, is carried out by hand.

\section{Consequences for the census}

The groups
\href{https://genus-2-torsion.github.io/group.html?g=30}{\([30]\)},
\href{https://genus-2-torsion.github.io/group.html?g=60}{\([60]\)} and
\href{https://genus-2-torsion.github.io/group.html?g=2.30}{\([2,30]\)}
are impossible in the class ``geometrically split, simple over
\(\Q\)''. They are the only census groups containing a point of order
\(30\).

Simplicity over \(\Q\) is essential. All three groups have examples in
the class ``split over \(\Q\)'':
\href{https://genus-2-torsion.github.io/curve.html?id=30.b}{30.b},
\href{https://genus-2-torsion.github.io/curve.html?id=60.a}{60.a} and
\href{https://genus-2-torsion.github.io/curve.html?id=2.30.a}{2.30.a}.
By Section~\ref{sec:hypotheses}, the two elliptic factors of each of
these Jacobians are not geometrically isogenous. Thus these exclusions
do not change the combined geometrically split ``\(\infty\)?'' grade of
any of the three groups, which remains certified exact, through
families split over \(\Q\). The group \([30]\) also keeps its
geometrically simple example
\href{https://genus-2-torsion.github.io/curve.html?id=30.a}{30.a}.

\begin{thebibliography}{9}
\bibitem{DN}
M. Derickx and F. Najman,
\emph{Classification of torsion of elliptic curves over quartic
fields}, arXiv:2412.16016v2 (2025), Theorem 1.1.
\href{https://arxiv.org/abs/2412.16016v2}{arXiv:2412.16016}.

\bibitem{Mazur}
B. Mazur,
\emph{Modular curves and the Eisenstein ideal}, Publications
Math\'ematiques de l'IH\'ES \textbf{47} (1977), 33--186, Theorem (8).
\href{https://www.numdam.org/item/PMIHES_1977__47__33_0/}{Numdam}.

\bibitem{KM}
M. A. Kenku and F. Momose,
\emph{Torsion points on elliptic curves defined over quadratic fields},
Nagoya Mathematical Journal \textbf{109} (1988), 125--149,
Theorem (0.1).
\href{https://doi.org/10.1017/S0027763000002816}{doi:10.1017/S0027763000002816}.

\bibitem{ST}
J.-P. Serre and J. Tate,
\emph{Good reduction of abelian varieties}, Annals of Mathematics
\textbf{88} (1968), 492--517: Theorem 1 and its Corollary 2 (p.~493),
Lemma 2 (p.~495), the case of a finite residue field in \S2 (p.~498),
and Theorem 6(a) (p.~503).
\href{https://doi.org/10.2307/1970722}{doi:10.2307/1970722}.

\bibitem{Tate}
J. Tate,
\emph{Endomorphisms of abelian varieties over finite fields},
Inventiones Mathematicae \textbf{2} (1966), 134--144, Theorem 2(c).
\href{https://doi.org/10.1007/BF01404549}{doi:10.1007/BF01404549}.

\bibitem{CX}
P. L. Clark and X. Xarles,
\emph{Local bounds for torsion points on abelian varieties},
Canadian Journal of Mathematics \textbf{60} (2008), 532--555, \S2.4
and Theorem 1.4. Theorem 1.4 concerns abelian surfaces with everywhere
anisotropic reduction, a class that includes potentially good
reduction.
\href{https://doi.org/10.4153/CJM-2008-026-x}{doi:10.4153/CJM-2008-026-x}.

\bibitem{Clark}
P. L. Clark,
\emph{Bounds for torsion on abelian varieties with integral moduli},
arXiv:math/0407264v2 (2004), \S2.1 and Theorem 5.
\href{https://arxiv.org/abs/math/0407264}{arXiv:math/0407264}.
\end{thebibliography}

\end{document}
